% =====================================================================
\chapter{Soluciones de los ejercicios}
\label{ap:soluciones}
% =====================================================================

Cada solución se presenta en cinco partes: datos, planteamiento, desarrollo, resultado y comprobación. En los pasos intermedios se redondea a cuatro cifras significativas salvo que se indique otra cosa; los logaritmos son naturales, y las divergencias se expresan en nats.

% ---------------------------------------------------------------------
\section{Tema 1}

\begin{solucion}{eje:1-chatbot}
\datos asistente que recibe una nota al final de cada respuesta.
\planteamiento aplicar la \cref{def:mdp-tokens} y la \cref{tab:mdp-llm}.
\pasos
\begin{enumerate}[label=\arabic*.]
  \item Estado $s_t$: la conversación previa (la instrucción $x$) más los tokens ya generados $y_{1:t}$; es todo lo que la política ve al decidir.
  \item Acción $a_t$: el siguiente token $y_{t+1}\in\V$.
  \item Transición: determinista; el nuevo estado es el anterior con el token concatenado.
  \item Recompensa: nula en cada token intermedio y igual a la nota $r(x,y)$ al emitir el símbolo de fin.
  \item Episodio: una respuesta completa, desde el primer token hasta el símbolo de fin o la longitud máxima.
  \item Con $\gamma=1$ y una sola recompensa final, el retorno es la nota (\cref{prop:mdp-tokens}), y la probabilidad de la trayectoria es $\pi_\theta(y\mid x)$; la respuesta entera puede tratarse como una sola acción en el contexto $x$.
\end{enumerate}
\resultado \textbf{estado $=$ contexto; acción $=$ token; transición $=$ concatenar; recompensa $=$ la nota, solo al final; episodio $=$ una respuesta}. Es un bandido contextual cuyas acciones son las respuestas.
\comprobacion la descripción coincide con la \cref{obs:bandido}: la dinámica es conocida y no hay recompensas intermedias que asignar.
\end{solucion}

\begin{solucion}{eje:1-contar}
\datos $\abs{\V}=5\cdot 10^{4}$; $T=100$.
\planteamiento hay $\abs{\V}^{T}$ respuestas; se toma $\log_{10}$.
\pasos
\begin{enumerate}[label=\arabic*.]
  \item $\log_{10}(5\cdot 10^{4})=\log_{10}5+4=0{,}69897+4=4{,}69897$.
  \item $\log_{10}\abs{\V}^{T}=100\cdot 4{,}69897=469{,}9$.
\end{enumerate}
\resultado \textbf{unas $10^{469{,}9}\approx 8\cdot 10^{469}$ respuestas}. La suma que define $V^{\pi}(x)$ no puede calcularse término a término; solo puede estimarse por muestreo, que es lo que hacen todos los métodos del módulo.
\comprobacion $10^{0{,}9}\approx 7{,}9$; y el número estimado de átomos del universo observable, del orden de $10^{80}$, es insignificante en comparación.
\end{solucion}

\begin{solucion}{eje:1-retornos}
\datos $\gamma=0{,}5$; $T=5$; $R_{1:5}=(-1,2,6,3,2)$.
\planteamiento recursión $G_t=R_{t+1}+\gamma G_{t+1}$ hacia atrás desde $G_5=0$ (\cref{prop:retorno}).
\pasos
\begin{enumerate}[label=\arabic*.]
  \item $G_5=0$.
  \item $G_4=R_5+0{,}5\,G_5=2+0=2$.
  \item $G_3=R_4+0{,}5\,G_4=3+1=4$.
  \item $G_2=R_3+0{,}5\,G_3=6+2=8$.
  \item $G_1=R_2+0{,}5\,G_2=2+4=6$.
  \item $G_0=R_1+0{,}5\,G_1=-1+3=2$.
\end{enumerate}
\resultado $\mathbf{(G_0,\dots,G_5)=(2,\,6,\,8,\,4,\,2,\,0)}$.
\comprobacion directamente, $G_0=-1+0{,}5\cdot 2+0{,}25\cdot 6+0{,}125\cdot 3+0{,}0625\cdot 2=-1+1+1{,}5+0{,}375+0{,}125=2$.
\end{solucion}

\begin{solucion}{eje:1-arbol}
\datos $\pi(\texttt{a}\mid x)=0{,}5$; $\pi(\texttt{a}\mid x\texttt{a})=0{,}2$; $\pi(\texttt{a}\mid x\texttt{b})=0{,}6$; $r(\texttt{aa})=0$, $r(\texttt{ab})=1$, $r(\texttt{ba})=2$, $r(\texttt{bb})=0$.
\planteamiento por la \cref{prop:mdp-tokens}, el valor de un prefijo es la recompensa esperada de sus continuaciones, y $Q(s,a)=V(sa)$.
\pasos
\begin{enumerate}[label=\arabic*.]
  \item $V(x\texttt{a})=0{,}2\cdot 0+0{,}8\cdot 1=0{,}8$.
  \item $V(x\texttt{b})=0{,}6\cdot 2+0{,}4\cdot 0=1{,}2$.
  \item $V(x)=0{,}5\cdot 0{,}8+0{,}5\cdot 1{,}2=1{,}0$.
  \item Primer token: $A(x,\texttt{a})=0{,}8-1{,}0=-0{,}2$; $A(x,\texttt{b})=1{,}2-1{,}0=0{,}2$.
  \item Segundo token tras \texttt{b}: $A(x\texttt{b},\texttt{a})=r(\texttt{ba})-V(x\texttt{b})=2-1{,}2=0{,}8$; $A(x\texttt{b},\texttt{b})=0-1{,}2=-1{,}2$.
\end{enumerate}
\resultado $\mathbf{V(x)=1}$, $\mathbf{V(x\texttt{a})=0{,}8}$, $\mathbf{V(x\texttt{b})=1{,}2}$; ventajas $(-0{,}2;\ 0{,}2)$ en el primer token y $(0{,}8;\ -1{,}2)$ en el segundo tras \texttt{b}.
\comprobacion \cref{cor:ventaja-media}: $0{,}5\cdot(-0{,}2)+0{,}5\cdot 0{,}2=0$ y $0{,}6\cdot 0{,}8+0{,}4\cdot(-1{,}2)=0{,}48-0{,}48=0$.
\end{solucion}

\begin{solucion}{eje:1-cs285}
\datos $\pi_\theta(a_1\mid s_1)=\theta$; $a_1$: recompensa~1 y permanece en $s_1$; $a_2$: recompensa~0 y termina; $V(s_F)=0$.
\planteamiento ecuación de Bellman (\cref{teo:bellman}) en el único estado no terminal.
\pasos
\begin{enumerate}[label=\arabic*.]
  \item $Q(s_1,a_1)=1+\gamma V(s_1)$ y $Q(s_1,a_2)=0+\gamma V(s_F)=0$.
  \item $V(s_1)=\theta\,(1+\gamma V(s_1))+(1-\theta)\cdot 0$, luego $V(s_1)(1-\gamma\theta)=\theta$ y $V(s_1)=\theta/(1-\gamma\theta)$.
  \item $\gamma=1$, $\theta=0{,}75$: $V=0{,}75/0{,}25=3$.
  \item $\gamma=0{,}9$, $\theta=0{,}75$: $V=0{,}75/(1-0{,}675)=0{,}75/0{,}325=2{,}308$.
  \item (c) El número $k$ de veces que se elige $a_1$ antes de $a_2$ cumple $P(k)=\theta^{k}(1-\theta)$, y $R(\tau)=k$. Entonces $\E[k]=(1-\theta)\sum_{k\ge0}k\theta^{k}=(1-\theta)\,\theta/(1-\theta)^{2}=\theta/(1-\theta)$, usando $\sum_k k\theta^{k}=\theta/(1-\theta)^{2}$ (derivada de la serie geométrica).
\end{enumerate}
\resultado $\mathbf{V(s_1)=\theta/(1-\gamma\theta)}$; vale $\mathbf{3}$ con $\gamma=1$ y $\mathbf{2{,}308}$ con $\gamma=0{,}9$.
\comprobacion con $\gamma=1$, $\E[k]=0{,}75/0{,}25=3$, igual que el paso~3. El descuento reduce el valor, como debe.
\end{solucion}

\begin{solucion}{eje:1-contraccion}
\datos $\gamma<1$; $(T^{\pi}v)(s)=\sum_a\pi(a\mid s)\sum_{s',r}p(s',r\mid s,a)[r+\gamma v(s')]$.
\planteamiento acotar la diferencia componente a componente; después usar que una contracción tiene a lo sumo un punto fijo, y expresar la ecuación en forma matricial.
\pasos
\begin{enumerate}[label=\arabic*.]
  \item Para cada $s$, los términos de recompensa se cancelan: $(T^{\pi}u)(s)-(T^{\pi}v)(s)=\gamma\sum_a\pi(a\mid s)\sum_{s',r}p(s',r\mid s,a)\,[u(s')-v(s')]$.
  \item Los pesos $\pi(a\mid s)p(s',r\mid s,a)$ son no negativos y suman~1, luego el valor absoluto del sumatorio es a lo sumo $\max_{s'}\abs{u(s')-v(s')}=\norm{u-v}_\infty$.
  \item Tomando el máximo en $s$: $\norm{T^{\pi}u-T^{\pi}v}_\infty\le\gamma\norm{u-v}_\infty$.
  \item (b) Si $u$ y $v$ son soluciones, $u=T^{\pi}u$ y $v=T^{\pi}v$, luego $\norm{u-v}_\infty\le\gamma\norm{u-v}_\infty$; como $\gamma<1$, $\norm{u-v}_\infty=0$.
  \item (c) Sean $P^{\pi}_{ss'}=\sum_a\pi(a\mid s)\sum_r p(s',r\mid s,a)$ y $r^{\pi}(s)=\sum_a\pi(a\mid s)\sum_{s',r}p(s',r\mid s,a)\,r$. La ecuación es $v=r^{\pi}+\gamma P^{\pi}v$, es decir, $(I-\gamma P^{\pi})v=r^{\pi}$.
  \item $P^{\pi}$ es estocástica por filas, luego $\norm{\gamma P^{\pi}w}_\infty\le\gamma\norm w_\infty$ para todo $w$. Si $(I-\gamma P^{\pi})w=0$, entonces $w=\gamma P^{\pi}w$ y $\norm w_\infty\le\gamma\norm w_\infty$, de donde $w=0$: la matriz tiene núcleo trivial y es invertible.
\end{enumerate}
\resultado \textbf{$T^{\pi}$ es una contracción de constante $\gamma$, y la ecuación de Bellman tiene solución única $v=(I-\gamma P^{\pi})^{-1}r^{\pi}$.}
\comprobacion en el \cref{eje:1-cs285}, con un único estado no terminal, $P^{\pi}=\theta$ y $r^{\pi}=\theta$, y la fórmula da $v=\theta/(1-\gamma\theta)$, el mismo resultado.
\end{solucion}

% ---------------------------------------------------------------------
\section{Tema 2}

\begin{solucion}{eje:2-baseline}
\planteamiento separar la esperanza del estimador (sesgo) de su dispersión (varianza).
\pasos
\begin{enumerate}[label=\arabic*.]
  \item El estimador con línea base es $(r-b)\,g$, con $g=\nabla_\theta\log\pi_\theta(a)$.
  \item Su esperanza es $\E[rg]-b\,\E[g]$, y $\E[g]=0$ por el \cref{lem:score-media}; luego $\E[(r-b)g]=\E[rg]=\nabla J$.
  \item La varianza depende de $\E[(r-b)^{2}\norm g^{2}]$, que sí cambia con $b$: si $b$ está cerca de las recompensas típicas, los factores $(r-b)$ son pequeños.
  \item Si $b$ dependiera de la acción, $\E[b(a)\,g(a)]=\sum_a\pi(a)b(a)\nabla\log\pi(a)=\sum_a b(a)\nabla\pi(a)$, que en general no es nulo, y el estimador quedaría sesgado.
\end{enumerate}
\resultado \textbf{la línea base no cambia la media porque la puntuación tiene media nula, pero sí el tamaño de cada muestra y, con ello, la varianza; debe ser independiente de la acción.}
\comprobacion en el \cref{ej:reinforce}, la media del estimador es $(0;\ 0{,}333;\ -0{,}333)$ con $b=0$ y con $b=1$, mientras que la varianza pasa de $0{,}889$ a $0{,}222$.
\end{solucion}

\begin{solucion}{eje:2-bernoulli}
\datos $\pi_\theta(1)=\sigma(\theta)$; $\theta=\ln 3$; $r(1)=1$, $r(0)=0$.
\planteamiento con $\sigma'=\sigma(1-\sigma)$: $\nabla\log\sigma(\theta)=1-\sigma(\theta)$ y $\nabla\log(1-\sigma(\theta))=-\sigma(\theta)$.
\pasos
\begin{enumerate}[label=\arabic*.]
  \item $e^{\theta}=3$, luego $\sigma(\ln 3)=3/(3+1)=0{,}75$.
  \item $\nabla_\theta\log\pi_\theta(1)=1-0{,}75=0{,}25$.
  \item $\nabla_\theta\log\pi_\theta(0)=-0{,}75$.
  \item $J(\theta)=\sigma(\theta)\cdot 1+(1-\sigma(\theta))\cdot 0=\sigma(\theta)$, luego $J'(\theta)=\sigma(1-\sigma)=0{,}75\cdot 0{,}25=0{,}1875$.
\end{enumerate}
\resultado $\mathbf{\pi_\theta(1)=0{,}75}$; puntuaciones $\mathbf{0{,}25}$ y $\mathbf{-0{,}75}$; $\mathbf{\nabla J=3/16=0{,}1875}$.
\comprobacion $\E[r\nabla\log\pi]=0{,}75\cdot 1\cdot 0{,}25+0{,}25\cdot 0\cdot(-0{,}75)=0{,}1875$. Además, $\E[\nabla\log\pi]=0{,}75\cdot 0{,}25+0{,}25\cdot(-0{,}75)=0$.
\end{solucion}

\begin{solucion}{eje:2-paso}
\datos $\theta=(0,0,0)$; $a=2$; $r=-1$; $\eta=1$; $b=0$.
\planteamiento $\theta\leftarrow\theta+\eta\,r\,(e_a-\pi)$ (\cref{prop:score-softmax,def:reinforce}).
\pasos
\begin{enumerate}[label=\arabic*.]
  \item $\pi=\softmax(0,0,0)=(1/3,1/3,1/3)$.
  \item $e_2-\pi=(-1/3,\,-1/3,\,2/3)$.
  \item $\theta=0+1\cdot(-1)\cdot(-1/3,-1/3,2/3)=(1/3,\,1/3,\,-2/3)$.
  \item $e^{1/3}=1{,}396$ y $e^{-2/3}=0{,}5134$; normalizador $1{,}396+1{,}396+0{,}5134=3{,}305$.
  \item $\pi=(1{,}396/3{,}305;\ 1{,}396/3{,}305;\ 0{,}5134/3{,}305)=(0{,}4223;\ 0{,}4223;\ 0{,}1554)$.
\end{enumerate}
\resultado $\mathbf{\theta=(1/3,1/3,-2/3)}$ y $\mathbf{\pi\approx(0{,}422;\ 0{,}422;\ 0{,}155)}$: la acción castigada pierde probabilidad.
\comprobacion $0{,}4223+0{,}4223+0{,}1554=1{,}000$, y por simetría las acciones~0 y~1 quedan iguales.
\end{solucion}

\begin{solucion}{eje:2-cs285}
\datos $p_\theta(\tau)=\theta^{k}(1-\theta)$; $R=k$; $\E[k]=\theta/(1-\theta)$; $\E[k^{2}]=\theta(1+\theta)/(1-\theta)^{2}$.
\planteamiento derivar $\log p_\theta(\tau)$ y aplicar el \cref{teo:pg}: $\nabla J=\E[R\,\nabla\log p_\theta(\tau)]$.
\pasos
\begin{enumerate}[label=\arabic*.]
  \item (a) $\log p_\theta(\tau)=k\log\theta+\log(1-\theta)$, luego $\nabla_\theta\log p_\theta(\tau)=k/\theta-1/(1-\theta)$.
  \item (b) $\nabla J=\E\bigl[k\,(k/\theta-1/(1-\theta))\bigr]=\E[k^{2}]/\theta-\E[k]/(1-\theta)$.
  \item $\E[k^{2}]/\theta=(1+\theta)/(1-\theta)^{2}$ y $\E[k]/(1-\theta)=\theta/(1-\theta)^{2}$.
  \item $\nabla J=\bigl((1+\theta)-\theta\bigr)/(1-\theta)^{2}=1/(1-\theta)^{2}$.
  \item (c) $J(\theta)=\theta/(1-\theta)$ (\cref{eje:1-cs285}); por la regla del cociente, $J'=\bigl((1-\theta)+\theta\bigr)/(1-\theta)^{2}=1/(1-\theta)^{2}$.
  \item En $\theta=1/2$: $1/(1/2)^{2}=4$.
\end{enumerate}
\resultado $\mathbf{\nabla_\theta\log p_\theta(\tau)=k/\theta-1/(1-\theta)}$ y $\mathbf{\nabla J=1/(1-\theta)^{2}}$, que vale $\mathbf{4}$ en $\theta=1/2$.
\comprobacion los dos caminos coinciden. Una simulación con $4\cdot 10^{6}$ trayectorias en $\theta=1/2$ da una media del estimador de $3{,}998$.
\end{solucion}

\begin{solucion}{eje:2-gae}
\datos $r=(1,0,2)$; $V=(2;\ 1{,}5;\ 1)$, $V(s_3)=0$; $\gamma=0{,}9$; $\lambda=0{,}8$, luego $\gamma\lambda=0{,}72$.
\planteamiento $\delta_t=r_{t+1}+\gamma V(s_{t+1})-V(s_t)$ y $\widehat A_t=\delta_t+\gamma\lambda\widehat A_{t+1}$ (\cref{def:gae}).
\pasos
\begin{enumerate}[label=\arabic*.]
  \item $\delta_0=1+0{,}9\cdot 1{,}5-2=0{,}35$.
  \item $\delta_1=0+0{,}9\cdot 1-1{,}5=-0{,}6$.
  \item $\delta_2=2+0-1=1$.
  \item $\widehat A_2=1$; $\widehat A_1=-0{,}6+0{,}72\cdot 1=0{,}12$; $\widehat A_0=0{,}35+0{,}72\cdot 0{,}12=0{,}35+0{,}0864=0{,}4364$.
  \item Con $\lambda=1$ ($\gamma\lambda=0{,}9$): $\widehat A_2=1$; $\widehat A_1=-0{,}6+0{,}9=0{,}3$; $\widehat A_0=0{,}35+0{,}27=0{,}62$.
  \item $G_0=1+0{,}9\cdot 0+0{,}81\cdot 2=2{,}62$, luego $G_0-V(s_0)=0{,}62$.
\end{enumerate}
\resultado $\mathbf{\delta=(0{,}35;\ -0{,}6;\ 1)}$ y $\mathbf{\widehat A=(0{,}4364;\ 0{,}12;\ 1)}$; con $\lambda=1$, $\widehat A_0=0{,}62=G_0-V(s_0)$.
\comprobacion la igualdad del paso~6 es la \cref{prop:gae}. El último paso no depende de $\lambda$ porque no hay errores posteriores.
\end{solucion}

\begin{solucion}{eje:2-optima}
\planteamiento primero la identidad de la puntuación; después, minimizar la varianza como función cuadrática de $b$.
\pasos
\begin{enumerate}[label=\arabic*.]
  \item $\E_{a\sim\pi}[\nabla\log\pi_\theta(a)]=\sum_a\pi_\theta(a)\,\nabla\pi_\theta(a)/\pi_\theta(a)=\sum_a\nabla\pi_\theta(a)=\nabla\sum_a\pi_\theta(a)=\nabla 1=0$.
  \item Sin sesgo: $\E[(r-b)g]=\E[rg]-b\,\E[g]=\E[rg]=\nabla J$, porque $b$ no depende de $a$.
  \item Varianza total: $V(b)=\E[\norm{(r-b)g}^{2}]-\norm{\E[(r-b)g]}^{2}=\E[(r-b)^{2}\norm g^{2}]-\norm{\nabla J}^{2}$; el segundo término no depende de $b$.
  \item $V'(b)=-2\,\E[(r-b)\norm g^{2}]=-2\,\E[r\norm g^{2}]+2b\,\E[\norm g^{2}]$.
  \item $V'(b)=0\iff b^{*}=\E[r\norm g^{2}]/\E[\norm g^{2}]$, y $V''(b)=2\E[\norm g^{2}]>0$, luego es un mínimo.
  \item Numéricamente, con $\pi=(0{,}5;\ 0{,}3;\ 0{,}2)$: $e_0-\pi=(0{,}5;\ -0{,}3;\ -0{,}2)$, de norma al cuadrado $0{,}25+0{,}09+0{,}04=0{,}38$; análogamente $0{,}78$ y $0{,}98$.
  \item Numerador: $0{,}5\cdot 1\cdot 0{,}38+0{,}3\cdot 2\cdot 0{,}78+0{,}2\cdot 0\cdot 0{,}98=0{,}19+0{,}468=0{,}658$. Denominador: $0{,}5\cdot 0{,}38+0{,}3\cdot 0{,}78+0{,}2\cdot 0{,}98=0{,}19+0{,}234+0{,}196=0{,}62$.
\end{enumerate}
\resultado $\mathbf{\E[\nabla\log\pi]=0}$, la línea base constante no introduce sesgo y $\mathbf{b^{*}=\E[r\norm g^{2}]/\E[\norm g^{2}]}$; en el ejemplo, $\mathbf{b^{*}=0{,}658/0{,}62=1{,}061}$.
\comprobacion $b^{*}$ es una media ponderada de las recompensas $(1,2,0)$ con pesos positivos, luego está entre $0$ y $2$; queda cerca de $\E[r]=1{,}1$.
\end{solucion}

\begin{solucion}{eje:2-varianza}
\datos $g(k)=k(k/\theta-1/(1-\theta))$; momentos $\E[k^{n}]$ del enunciado y del \cref{eje:2-cs285}; $\E[g]=1/(1-\theta)^{2}$.
\planteamiento $\Var[g]=\E[g^{2}]-\E[g]^{2}$, con $g^{2}$ desarrollado como polinomio en $k$.
\pasos
\begin{enumerate}[label=\arabic*.]
  \item $g^{2}=k^{2}\bigl(k^{2}/\theta^{2}-2k/(\theta(1-\theta))+1/(1-\theta)^{2}\bigr)=k^{4}/\theta^{2}-2k^{3}/(\theta(1-\theta))+k^{2}/(1-\theta)^{2}$.
  \item Se escriben los tres términos con denominador común $\theta(1-\theta)^{4}$:
  \[
  \frac{\E[k^{4}]}{\theta^{2}}=\frac{(1+\theta)(\theta^{2}+10\theta+1)}{\theta(1-\theta)^{4}},\quad \frac{2\E[k^{3}]}{\theta(1-\theta)}=\frac{2\theta(\theta^{2}+4\theta+1)}{\theta(1-\theta)^{4}},\quad \frac{\E[k^{2}]}{(1-\theta)^{2}}=\frac{\theta^{2}(1+\theta)}{\theta(1-\theta)^{4}}.
  \]
  \item Numerador: $(\theta^{3}+11\theta^{2}+11\theta+1)-(2\theta^{3}+8\theta^{2}+2\theta)+(\theta^{3}+\theta^{2})=4\theta^{2}+9\theta+1$.
  \item $\E[g]^{2}=1/(1-\theta)^{4}=\theta/\bigl(\theta(1-\theta)^{4}\bigr)$.
  \item $\Var[g]=(4\theta^{2}+9\theta+1-\theta)/\bigl(\theta(1-\theta)^{4}\bigr)=(4\theta^{2}+8\theta+1)/\bigl(\theta(1-\theta)^{4}\bigr)$.
  \item En $\theta=1/2$: numerador $1+4+1=6$; denominador $\tfrac12\cdot\tfrac1{16}=\tfrac1{32}$; $\Var=6\cdot 32=192$.
  \item Cuando $\theta\to 0$, $\Var\sim 1/\theta\to\infty$; cuando $\theta\to 1$, $\Var\sim 13/(1-\theta)^{4}\to\infty$.
\end{enumerate}
\resultado $\mathbf{\Var[g]=(4\theta^{2}+8\theta+1)/(\theta(1-\theta)^{4})}$, que vale $\mathbf{192}$ en $\theta=1/2$ y diverge en ambos extremos; su mínimo, hallado numéricamente, está en $\theta\approx 0{,}110$, con $\Var\approx 27{,}9$.
\comprobacion en $\theta=1/2$, $\E[g^{2}]=(1+4{,}5+1)\cdot 32=208$ y $208-4^{2}=192$. La simulación de $4\cdot 10^{6}$ trayectorias da una varianza muestral de $193$. La desviación típica, $\sqrt{192}\approx 13{,}9$, es más de tres veces el gradiente que se quiere estimar: de ahí la necesidad de líneas base.
\end{solucion}

% ---------------------------------------------------------------------
\section{Tema 3}

\begin{solucion}{eje:3-pesimista}
\planteamiento aplicar la \cref{prop:ppo-casos} y la discusión de la penalización KL.
\pasos
\begin{enumerate}[label=(\alph*)]
  \item \textbf{Falsa.} $\min(\rho A,\clip(\rho)A)\le\rho A$: el recorte solo puede bajar o dejar igual el objetivo.
  \item \textbf{Falsa.} Con $A>0$, $\ell(\rho)=\min(\rho,1+\varepsilon)A=\rho A$ si $\rho<1-\varepsilon$; la derivada es $A\neq 0$. El gradiente se anula con $A>0$ solo si $\rho>1+\varepsilon$.
  \item \textbf{Verdadera.} $\E_\pi[r]=\sum_y\pi(y)r(y)$ es lineal en $\pi$ y su máximo sobre el símplex se alcanza en un vértice, la masa puntual en una respuesta de máxima recompensa (si hay varias empatadas, cualquier mezcla de ellas es óptima).
\end{enumerate}
\resultado \textbf{(a) falsa; (b) falsa; (c) verdadera.}
\comprobacion en la \cref{tab:ppo} el objetivo es siempre menor o igual que $\rho A$ ($2{,}4\le3{,}0$; $1{,}0\le1{,}0$; $-1{,}5\le-1{,}5$; $-0{,}8\le-0{,}5$), y la fila $(0{,}5;\ +2)$ conserva el gradiente.
\end{solucion}

\begin{solucion}{eje:3-ppo}
\datos $\varepsilon=0{,}2$, luego $\clip(\rho)\in[0{,}8;\ 1{,}2]$.
\planteamiento calcular $\rho A$ y $\clip(\rho)A$ y tomar el mínimo; si gana el término recortado con $\rho$ fuera del intervalo, el gradiente es nulo.
\pasos
\begin{enumerate}[label=\arabic*.]
  \item $(1{,}3;\ 1)$: $\rho A=1{,}3$; $\clip(\rho)A=1{,}2$; mínimo $1{,}2$, término recortado: gradiente nulo.
  \item $(0{,}7;\ 1)$: $\rho A=0{,}7$; $\clip(\rho)A=0{,}8$; mínimo $0{,}7$: gradiente no nulo.
  \item $(0{,}7;\ -2)$: $\rho A=-1{,}4$; $\clip(\rho)A=0{,}8\cdot(-2)=-1{,}6$; mínimo $-1{,}6$, recortado: gradiente nulo.
  \item $(1{,}3;\ -2)$: $\rho A=-2{,}6$; $\clip(\rho)A=1{,}2\cdot(-2)=-2{,}4$; mínimo $-2{,}6$: gradiente no nulo.
\end{enumerate}
\resultado \textbf{objetivos $1{,}2$; $0{,}7$; $-1{,}6$; $-2{,}6$. Gradiente nulo en el primero y en el tercero.}
\comprobacion el gradiente se anula cuando el cociente ya se movió más de $\varepsilon$ en la dirección que favorece la ventaja (subir si $A>0$, bajar si $A<0$), como predice la \cref{prop:ppo-casos}.
\end{solucion}

\begin{solucion}{eje:3-grpoclip}
\datos $\varepsilon=0{,}2$; $A=+1$ con $\Delta=0{,}1;\ 0{,}3;\ -0{,}4$; $A=-1$ con $\Delta=0{,}3;\ -0{,}4$; $\rho=e^{\Delta}$.
\planteamiento calcular $\rho$, recortarlo a $[0{,}8;\ 1{,}2]$, tomar el mínimo y cambiar el signo.
\pasos
\begin{enumerate}[label=\arabic*.]
  \item $e^{0{,}1}=1{,}105$; $e^{0{,}3}=1{,}350$; $e^{-0{,}4}=0{,}6703$.
  \item $A=+1$, $\rho=1{,}105$ (dentro del intervalo): mínimo $1{,}105$; pérdida $-1{,}105$; gradiente no nulo.
  \item $A=+1$, $\rho=1{,}350$: $\min(1{,}350;\ 1{,}2)=1{,}2$; pérdida $-1{,}2$; gradiente nulo.
  \item $A=+1$, $\rho=0{,}6703$: $\min(0{,}6703;\ 0{,}8)=0{,}6703$; pérdida $-0{,}6703$; gradiente no nulo.
  \item $A=-1$, $\rho=1{,}350$: $\min(-1{,}350;\ -1{,}2)=-1{,}350$; pérdida $+1{,}350$; gradiente no nulo.
  \item $A=-1$, $\rho=0{,}6703$: $\min(-0{,}6703;\ -0{,}8)=-0{,}8$; pérdida $+0{,}8$; gradiente nulo.
\end{enumerate}
\resultado \textbf{pérdidas $-1{,}105$; $-1{,}2$; $-0{,}670$; $+1{,}350$; $+0{,}8$. Gradiente nulo en el segundo ($A>0$, $\rho>1{,}2$) y en el quinto ($A<0$, $\rho<0{,}8$).}
\comprobacion el tercer y el cuarto token (pasos~4 y~5) son los que se movieron en la dirección desfavorable y conservan el gradiente, igual que en la \cref{tab:ppo}.
\end{solucion}

\begin{solucion}{eje:3-beta}
\datos $\piref=(\tfrac12,\tfrac12)$; $r=(1,0)$; $\beta\in\{1;\ 0{,}5\}$.
\planteamiento $\pi^{*}(y)\propto\piref(y)e^{r(y)/\beta}$ (\cref{teo:optimo-kl}); después $\KL=\sum\pi^{*}\log(\pi^{*}/\piref)$.
\pasos
\begin{enumerate}[label=\arabic*.]
  \item $\pi^{*}(0)=\tfrac12e^{1/\beta}/(\tfrac12e^{1/\beta}+\tfrac12)=e^{1/\beta}/(e^{1/\beta}+1)=\sigma(1/\beta)$.
  \item $\beta=1$: $e/(e+1)=2{,}718/3{,}718=0{,}7311$.
  \item $\beta=0{,}5$: $e^{2}/(e^{2}+1)=7{,}389/8{,}389=0{,}8808$.
  \item KL con $\beta=1$: $0{,}7311\ln(0{,}7311/0{,}5)=0{,}7311\cdot 0{,}3799=0{,}2777$.
  \item $0{,}2689\ln(0{,}2689/0{,}5)=0{,}2689\cdot(-0{,}6201)=-0{,}1668$.
  \item $\KL=0{,}2777-0{,}1668=0{,}1109$.
\end{enumerate}
\resultado $\mathbf{\pi^{*}=(0{,}731;\ 0{,}269)}$ con $\beta=1$ y $\mathbf{(0{,}881;\ 0{,}119)}$ con $\beta=0{,}5$; $\mathbf{\KL=0{,}111}$ nats.
\comprobacion un $\beta$ menor concentra más la política en la mejor acción, como en la \cref{tab:beta}; además $\KL\ge 0$.
\end{solucion}

\begin{solucion}{eje:3-estimadores}
\datos $\pi_\theta=(0{,}6;\ 0{,}3;\ 0{,}1)$; $\piref=(0{,}4;\ 0{,}4;\ 0{,}2)$; $\KL=0{,}0877$ (\cref{ej:kl}).
\planteamiento evaluar $k_1$ y $k_3$ en cada token y promediar con los pesos $\pi_\theta$.
\pasos
\begin{enumerate}[label=\arabic*.]
  \item $k_1$: $\ln(0{,}6/0{,}4)=0{,}4055$; $\ln(0{,}3/0{,}4)=-0{,}2877$; $\ln(0{,}1/0{,}2)=-0{,}6931$.
  \item $\E[k_1]=0{,}6\cdot 0{,}4055+0{,}3\cdot(-0{,}2877)+0{,}1\cdot(-0{,}6931)=0{,}2433-0{,}0863-0{,}0693=0{,}0877$.
  \item $k_3$, token~1: $u=0{,}4/0{,}6=0{,}6667$; $0{,}6667-\ln 0{,}6667-1=0{,}6667+0{,}4055-1=0{,}0721$.
  \item Token~2: $u=0{,}4/0{,}3=1{,}333$; $1{,}333-0{,}2877-1=0{,}0457$.
  \item Token~3: $u=2$; $2-0{,}6931-1=0{,}3069$.
  \item $\E[k_3]=0{,}6\cdot 0{,}0721+0{,}3\cdot 0{,}0457+0{,}1\cdot 0{,}3069=0{,}0433+0{,}0137+0{,}0307=0{,}0877$.
\end{enumerate}
\resultado $\mathbf{k_1=(0{,}405;\ -0{,}288;\ -0{,}693)}$ y $\mathbf{k_3=(0{,}072;\ 0{,}046;\ 0{,}307)}$; ambos tienen esperanza $\mathbf{0{,}0877}$ nats.
\comprobacion $k_1$ es negativo en dos tokens, de modo que una penalización por token basada en $k_1$ puede actuar como premio en una muestra concreta; $k_3$ es siempre no negativo (\cref{prop:k3}).
\end{solucion}

\begin{solucion}{eje:3-limites}
\datos $\beta\log Z$, con $Z=\sum_y\piref(y)e^{r(y)/\beta}$; sea $r_{\max}=\max\{r(y):\piref(y)>0\}$.
\planteamiento sacar factor común $e^{r_{\max}/\beta}$ para $\beta\to 0$; desarrollar en serie para $\beta\to\infty$.
\pasos
\begin{enumerate}[label=\arabic*.]
  \item $Z=e^{r_{\max}/\beta}\,S(\beta)$ con $S(\beta)=\sum_y\piref(y)e^{(r(y)-r_{\max})/\beta}$, luego $\beta\log Z=r_{\max}+\beta\log S(\beta)$.
  \item Cuando $\beta\to 0^{+}$, los términos con $r(y)<r_{\max}$ tienden a~0 y los de $r(y)=r_{\max}$ valen $\piref(y)$, de modo que $S(\beta)\to c\in(0,1]$ y $\beta\log S(\beta)\to 0$. Límite: $r_{\max}$.
  \item Cuando $\beta\to\infty$, $e^{r/\beta}=1+r/\beta+O(\beta^{-2})$, luego $Z=1+\E_{\piref}[r]/\beta+O(\beta^{-2})$.
  \item $\log(1+u)=u+O(u^{2})$, luego $\beta\log Z=\E_{\piref}[r]+O(1/\beta)\to\E_{\piref}[r]$.
  \item Con los datos del \cref{ej:beta}: $\beta=0{,}01$ da $\beta\log Z=1{,}988$, y $\beta=100$ da $1{,}102$.
\end{enumerate}
\resultado \textbf{el valor óptimo tiende a $r_{\max}$ cuando $\beta\to 0$ y a $\E_{\piref}[r]$ cuando $\beta\to\infty$.}
\comprobacion en el ejemplo, $r_{\max}=2$ y $\E_{\piref}[r]=0{,}5+0{,}6+0=1{,}1$, a los que se aproximan $1{,}988$ y $1{,}102$. Para $\beta=0{,}01$, $1{,}988=2+0{,}01\ln 0{,}3$, de acuerdo con el paso~1 ($S\to\piref(y=1)=0{,}3$).
\end{solucion}

% ---------------------------------------------------------------------
\section{Tema 4}

\begin{solucion}{eje:4-conceptos}
\planteamiento recordar la \cref{prop:bt} y la motivación del tema.
\pasos
\begin{enumerate}[label=(\alph*)]
  \item Las notas absolutas varían mucho entre anotadores y, para un mismo anotador, con el tiempo y el contexto; elegir entre dos respuestas es una tarea más fácil y más consistente, y el modelo de Bradley--Terry convierte esas elecciones en puntuaciones.
  \item \textbf{Verdadera.} $\sigma\bigl((r_a+5)-(r_b+5)\bigr)=\sigma(r_a-r_b)$.
  \item Multiplicar por~2 sí cambia las probabilidades: $\sigma(2(r_a-r_b))$ es más extrema que $\sigma(r_a-r_b)$. La escala de las recompensas está determinada por los datos; solo la traslación es arbitraria.
\end{enumerate}
\resultado \textbf{(a) por consistencia y facilidad; (b) verdadera; (c) cambia las probabilidades, porque la escala sí está identificada.}
\comprobacion con $r_a-r_b=1$, $\sigma(1)=0{,}731$ y $\sigma(2)=0{,}881$.
\end{solucion}

\begin{solucion}{eje:4-bt}
\datos $z=r_w-r_l=1$.
\planteamiento $P=\sigma(z)$; pérdida $-\ln\sigma(z)$; derivada $-(1-\sigma(z))$ (\cref{prop:bt-mv}).
\pasos
\begin{enumerate}[label=\arabic*.]
  \item $e^{-1}=0{,}3679$; $\sigma(1)=1/1{,}3679=0{,}7311$.
  \item Pérdida: $-\ln 0{,}7311=0{,}3133$.
  \item Derivada respecto a $r_w$: $-(1-0{,}7311)=-0{,}2689$.
\end{enumerate}
\resultado $\mathbf{P=0{,}731}$; pérdida $\mathbf{0{,}313}$; $\mathbf{\partial L/\partial r_w=-0{,}269}$.
\comprobacion la derivada es negativa, de modo que el descenso de gradiente sube $r_w$; su módulo es menor que el $0{,}5$ del caso $r_w=r_l$, porque el par ya está bien ordenado.
\end{solucion}

\begin{solucion}{eje:4-cs234}
\datos $\sigma^{1}$ con recompensas $(0{,}5;\ 1{,}0)$; $\sigma^{2}$ con $(0{,}2;\ 0{,}3)$; $\mu\in\{1;\ 0{,}5\}$.
\planteamiento $\frac{d}{dz}\log\sigma(z)=1-\sigma(z)$ y $\frac{d}{dz}\log(1-\sigma(z))=-\sigma(z)$; después, regla de la cadena con $z(\theta)$.
\pasos
\begin{enumerate}[label=\arabic*.]
  \item $dL/dz=-[\mu(1-\sigma(z))-(1-\mu)\sigma(z)]=-(\mu-\sigma(z))$.
  \item $\nabla_\theta z=\sum_t\nabla_\theta\hat r_\theta(o^{1}_t,a^{1}_t)-\sum_t\nabla_\theta\hat r_\theta(o^{2}_t,a^{2}_t)$.
  \item (a) $\nabla_\theta L=-(\mu-\sigma(z))\bigl[\sum_t\nabla_\theta\hat r_\theta(o^{1}_t,a^{1}_t)-\sum_t\nabla_\theta\hat r_\theta(o^{2}_t,a^{2}_t)\bigr]$.
  \item (b) $R^{1}=1{,}5$; $R^{2}=0{,}5$; $z=1$; $P=\sigma(1)=0{,}7311$.
  \item $\mu=1$: $L=-\ln 0{,}7311=0{,}3133$; $dL/dz=-(1-0{,}7311)=-0{,}2689$.
  \item $\mu=0{,}5$: $\ln 0{,}7311=-0{,}3133$ y $\ln 0{,}2689=-1{,}313$; $L=0{,}5\cdot 0{,}3133+0{,}5\cdot 1{,}313=0{,}8133$; $dL/dz=-(0{,}5-0{,}7311)=+0{,}2311$.
\end{enumerate}
\resultado $\mathbf{\nabla L=-(\mu-\sigma(z))(\sum\nabla\hat r^{1}-\sum\nabla\hat r^{2})}$; $\mathbf{P=0{,}731}$; con $\mu=1$, $L=0{,}313$ y $dL/dz=-0{,}269$; con $\mu=0{,}5$, $L=0{,}813$ y $dL/dz=+0{,}231$.
\comprobacion con $\mu=0{,}5$ el gradiente es positivo: el descenso reduce $z$ hacia~0, donde $P=0{,}5$ coincide con la etiqueta de empate; allí la pérdida sería $\ln 2=0{,}693<0{,}813$.
\end{solucion}

\begin{solucion}{eje:4-ranking}
\datos $K\in\{4,9\}$; ordenación $A\pref B\pref C$ con $r=(1{,}0;\ 0{,}2;\ -0{,}5)$.
\planteamiento una ordenación de $K$ elementos determina el resultado de cada uno de los $\binom K2$ pares.
\pasos
\begin{enumerate}[label=\arabic*.]
  \item $\binom 42=6$ y $\binom 92=36$.
  \item Pares de la ordenación: $(A,B)$ con diferencia $0{,}8$; $(A,C)$ con $1{,}5$; $(B,C)$ con $0{,}7$.
  \item $-\ln\sigma(0{,}8)=0{,}3711$; $-\ln\sigma(1{,}5)=0{,}2014$; $-\ln\sigma(0{,}7)=0{,}4032$.
  \item Media ponderada por $1/3$: $(0{,}3711+0{,}2014+0{,}4032)/3=0{,}9757/3=0{,}3252$.
\end{enumerate}
\resultado \textbf{6 y 36 comparaciones; pérdidas $0{,}371$, $0{,}201$ y $0{,}403$; media $0{,}325$.}
\comprobacion la pérdida más baja corresponde al par con mayor margen ($A,C$), y todas son menores que $\ln 2=0{,}693$ porque el modelo ordena bien los tres pares.
\end{solucion}

\begin{solucion}{eje:4-mv}
\datos $n=10$; $n_a=7$; $b\pref c$ con razón $2{:}1$.
\planteamiento \cref{prop:bt-dos} y \cref{prop:bt}(iii).
\pasos
\begin{enumerate}[label=\arabic*.]
  \item $\widehat{r_a-r_b}=\ln(7/3)=0{,}8473$.
  \item $\sigma(0{,}8473)=1/(1+3/7)=7/10=0{,}7$.
  \item Razón $a{:}b$ estimada: $7/3$. Razón $a{:}c$: $(7/3)\cdot 2=14/3$.
  \item $P(a\pref c)=(14/3)/(1+14/3)=14/17=0{,}8235$.
\end{enumerate}
\resultado la diferencia estimada es $\mathbf{0{,}847}$ y reproduce la frecuencia observada $\mathbf{0{,}7}$; el modelo predice $\mathbf{P(a\pref c)=14/17\approx 0{,}824}$.
\comprobacion $r_a-r_c=\ln(7/3)+\ln 2=\ln(14/3)=1{,}540$ y $\sigma(1{,}540)=0{,}824$.
\end{solucion}

\begin{solucion}{eje:4-convexidad}
\datos $r_u(x,y)=u^{\top}\varphi(x,y)$; $\Delta_m=\varphi(x^{(m)},y_w^{(m)})-\varphi(x^{(m)},y_l^{(m)})$; $z_m=u^{\top}\Delta_m$.
\planteamiento calcular gradiente y hessiana de $\LRM(u)=-\frac1M\sum_m\log\sigma(u^{\top}\Delta_m)$.
\pasos
\begin{enumerate}[label=\arabic*.]
  \item Gradiente: $\nabla\LRM=-\frac1M\sum_m(1-\sigma(z_m))\,\Delta_m$.
  \item Como $\frac{d}{dz}(1-\sigma(z))=-\sigma(z)(1-\sigma(z))$, la hessiana es $H=\frac1M\sum_m\sigma(z_m)(1-\sigma(z_m))\,\Delta_m\Delta_m^{\top}$.
  \item Para todo $v$, $v^{\top}Hv=\frac1M\sum_m\sigma(z_m)(1-\sigma(z_m))(v^{\top}\Delta_m)^{2}\ge 0$: $H$ es semidefinida positiva y $\LRM$ es convexa.
  \item (b) Cada término $-\log\sigma(z)$ es positivo, luego $\LRM>0$. Si $u_0^{\top}\Delta_m>0$ para todo $m$, con $u=cu_0$ y $c\to\infty$ cada $z_m\to\infty$ y $-\log\sigma(z_m)\to 0$: el ínfimo es~0, pero ningún $u$ finito lo alcanza.
  \item (c) Si los datos son separables, la máxima verosimilitud empuja la norma de $u$ hacia infinito y las probabilidades hacia 0 o 1: el modelo se vuelve arbitrariamente confiado, un caso extremo de sobreajuste que la regularización ($\ell_2$) o la parada temprana evitan.
\end{enumerate}
\resultado \textbf{$H=\frac1M\sum_m\sigma(z_m)(1-\sigma(z_m))\Delta_m\Delta_m^{\top}\succeq 0$; con datos separables, $\inf\LRM=0$ no se alcanza.}
\comprobacion es la misma estructura que la regresión logística del módulo~8, con características $\Delta_m$ y todas las etiquetas iguales a~1. La convexidad se pierde si $r_\phi$ es una red neuronal.
\end{solucion}

% ---------------------------------------------------------------------
\section{Tema 5}

\begin{solucion}{eje:5-vf}
\planteamiento repasar la \cref{def:dpo}, la \cref{prop:pref-pol} y la \cref{tab:dpo-ppo}.
\pasos
\begin{enumerate}[label=(\alph*)]
  \item \textbf{Falsa.} DPO usa la recompensa implícita $\hat r_\theta=\beta\log(\pi_\theta/\piref)$; no entrena ninguna red de recompensa aparte.
  \item \textbf{Falsa.} Trabaja con comparaciones $(x,y_w,y_l)$ recogidas de antemano.
  \item \textbf{Falsa.} $Z(x)$ no vale~1 en general; se cancela porque depende solo de $x$ y las dos respuestas comparten $x$.
  \item \textbf{Verdadera.} Con $\pi_\theta=\piref$, ambas recompensas implícitas son nulas, $h=0$ y $-\log\sigma(0)=\ln 2$.
\end{enumerate}
\resultado \textbf{(a) F, (b) F, (c) F, (d) V.}
\comprobacion $\ln 2=0{,}693$ es el valor inicial citado en el \cref{ej:dpo}.
\end{solucion}

\begin{solucion}{eje:5-paso}
\datos $\beta=0{,}5$; diferencias de log-probabilidad $+0{,}4$ (preferida) y $-0{,}6$ (rechazada).
\planteamiento $\hat r=\beta(\log\pi_\theta-\log\piref)$; $h=\hat r_w-\hat r_l$; pérdida $-\ln\sigma(h)$; peso $\beta\sigma(-h)$.
\pasos
\begin{enumerate}[label=\arabic*.]
  \item $\hat r_w=0{,}5\cdot 0{,}4=0{,}2$; $\hat r_l=0{,}5\cdot(-0{,}6)=-0{,}3$.
  \item $h=0{,}2-(-0{,}3)=0{,}5$.
  \item $e^{-0{,}5}=0{,}6065$; $\sigma(0{,}5)=1/1{,}6065=0{,}6225$; pérdida $-\ln 0{,}6225=0{,}4741$.
  \item $\sigma(-0{,}5)=0{,}3775$; peso $0{,}5\cdot 0{,}3775=0{,}1888$.
\end{enumerate}
\resultado $\mathbf{\hat r_w=0{,}2}$; $\mathbf{\hat r_l=-0{,}3}$; $\mathbf{h=0{,}5}$; pérdida $\mathbf{0{,}474}$; peso $\mathbf{0{,}189}$.
\comprobacion la pérdida queda por debajo de $\ln 2$ y el peso por debajo de $\beta/2=0{,}25$, porque el par ya está bien ordenado ($h>0$).
\end{solucion}

\begin{solucion}{eje:5-cs336}
\datos $\beta=0{,}2$; log-probabilidades por token del enunciado.
\planteamiento la log-probabilidad de una secuencia es la suma de las de sus tokens; después, la \cref{def:dpo}.
\pasos
\begin{enumerate}[label=\arabic*.]
  \item $\log\pi_\theta(y_w)=-0{,}5-1{,}2-0{,}3=-2{,}0$; $\log\piref(y_w)=-0{,}6-1{,}5-0{,}4=-2{,}5$.
  \item $\log\pi_\theta(y_l)=-1{,}0-2{,}0=-3{,}0$; $\log\piref(y_l)=-0{,}8-1{,}7=-2{,}5$.
  \item $\hat r_w=0{,}2\cdot(-2{,}0+2{,}5)=0{,}1$; $\hat r_l=0{,}2\cdot(-3{,}0+2{,}5)=-0{,}1$.
  \item $h=0{,}1-(-0{,}1)=0{,}2$.
  \item $e^{-0{,}2}=0{,}8187$; $\sigma(0{,}2)=1/1{,}8187=0{,}5498$; pérdida $-\ln 0{,}5498=0{,}5981$.
  \item $\sigma(-0{,}2)=0{,}4502$; peso $0{,}2\cdot 0{,}4502=0{,}0900$.
\end{enumerate}
\resultado $\mathbf{\hat r_w=0{,}1}$; $\mathbf{\hat r_l=-0{,}1}$; $\mathbf{h=0{,}2}$; pérdida $\mathbf{0{,}598}$; peso $\mathbf{0{,}0900}$.
\comprobacion $0{,}598<\ln 2$ porque $h>0$, y el peso queda algo por debajo de $\beta/2=0{,}1$. Las respuestas tienen distinta longitud, pero DPO solo usa las sumas.
\end{solucion}

\begin{solucion}{eje:5-tabular}
\datos $\piref=(\tfrac12,\tfrac12)$; $\beta=0{,}5$; $P(y_1\pref y_2)=p=0{,}8$ en los datos.
\planteamiento promediar la pérdida sobre los dos órdenes posibles, minimizar en $h$ y despejar $\pi_\theta$.
\pasos
\begin{enumerate}[label=\arabic*.]
  \item Con probabilidad $p$ el par es $(y_w,y_l)=(y_1,y_2)$ y su margen es $h$; con probabilidad $1-p$ es $(y_2,y_1)$ y su margen es $-h$. La pérdida esperada es $\ell(h)=-p\log\sigma(h)-(1-p)\log\sigma(-h)$.
  \item $\ell'(h)=-p\,\sigma(-h)+(1-p)\,\sigma(h)=-p+\sigma(h)$, usando $\sigma(-h)=1-\sigma(h)$; se anula en $\sigma(h)=p$, y $\ell''=\sigma(h)(1-\sigma(h))>0$.
  \item $h^{*}=\log\bigl(p/(1-p)\bigr)=\ln 4=1{,}386$.
  \item $h=\beta\bigl[\log(\pi(y_1)/\piref(y_1))-\log(\pi(y_2)/\piref(y_2))\bigr]=\beta\log(\pi(y_1)/\pi(y_2))$, porque $\piref$ es uniforme.
  \item $\log(\pi(y_1)/\pi(y_2))=1{,}386/0{,}5=2{,}773=\ln 16$, luego $\pi(y_1)/\pi(y_2)=16$ y $\pi_\theta=(16/17,\ 1/17)=(0{,}941;\ 0{,}059)$.
  \item Con $r=(\ln 4,0)$ y $\beta=0{,}5$: $\pi^{*}\propto(\tfrac12e^{2\ln 4},\ \tfrac12)=(8;\ 0{,}5)$, es decir, $\pi^{*}=(16/17,\ 1/17)$.
\end{enumerate}
\resultado $\mathbf{h^{*}=\ln 4}$ y $\mathbf{\pi_\theta=(16/17,\ 1/17)\approx(0{,}941;\ 0{,}059)}$, que coincide con $\pi^{*}$ para $r=(\ln 4,0)$.
\comprobacion bajo Bradley--Terry, $r=(\ln 4,0)$ da $P(y_1\pref y_2)=\sigma(\ln 4)=4/5=0{,}8=p$: DPO recupera exactamente la política óptima de la recompensa que generó los datos, como afirma el \cref{teo:representacion}.
\end{solucion}

\begin{solucion}{eje:5-gradiente}
\planteamiento regla de la cadena en dos eslabones: $L=-\log\sigma(h)$ y $h=h(\theta)$.
\pasos
\begin{enumerate}[label=\arabic*.]
  \item $dL/dh=-\sigma'(h)/\sigma(h)=-\sigma(h)(1-\sigma(h))/\sigma(h)=-(1-\sigma(h))=-\sigma(-h)$.
  \item $h=\beta[\log\pi_\theta(y_w\mid x)-\log\piref(y_w\mid x)]-\beta[\log\pi_\theta(y_l\mid x)-\log\piref(y_l\mid x)]$.
  \item $\piref$ está congelado: $\nabla_\theta\log\piref=0$.
  \item $\nabla_\theta h=\beta\bigl(\nabla_\theta\log\pi_\theta(y_w\mid x)-\nabla_\theta\log\pi_\theta(y_l\mid x)\bigr)$.
  \item $\nabla_\theta L=(dL/dh)\,\nabla_\theta h=-\beta\,\sigma(-h)\bigl(\nabla_\theta\log\pi_\theta(y_w\mid x)-\nabla_\theta\log\pi_\theta(y_l\mid x)\bigr)$; la esperanza sobre $\mathcal{D}$ da la \cref{prop:grad-dpo}.
  \item En $\pi_\theta=\piref$, $h=0$ y $\beta\sigma(0)=\beta/2$.
\end{enumerate}
\resultado $\mathbf{\nabla\LDPO=-\beta\,\E[\sigma(-h)(\nabla\log\pi_\theta(y_w\mid x)-\nabla\log\pi_\theta(y_l\mid x))]}$, con peso inicial $\beta/2$.
\comprobacion en el \cref{eje:5-paso} ($h=0{,}5$) el peso baja a $0{,}189<0{,}25=\beta/2$, como corresponde a un par ya bien ordenado.
\end{solucion}

\begin{solucion}{eje:5-invariancia}
\datos $r'(x,y)=r(x,y)+f(x)$.
\planteamiento sustituir $r'$ en $\pi^{*}$ y en Bradley--Terry; comparar con la reparametrización del \cref{lem:reparam}.
\pasos
\begin{enumerate}[label=\arabic*.]
  \item $\piref\,e^{r'/\beta}=\piref\,e^{r/\beta}\,e^{f(x)/\beta}$, y $Z'(x)=Z(x)\,e^{f(x)/\beta}$; al normalizar, el factor se cancela y $\pi^{*}_{r'}=\pi^{*}_r$ (\cref{lem:equiv-pol}).
  \item $r'(x,y_w)-r'(x,y_l)=r(x,y_w)-r(x,y_l)$, luego las probabilidades y la pérdida de Bradley--Terry no cambian (\cref{lem:equiv-pref}).
  \item En DPO, $r=\beta\log(\pi^{*}/\piref)+\beta\log Z(x)$: la recompensa es la implícita más un término $f(x)=\beta\log Z(x)$.
  \item Por el paso~2, Bradley--Terry no distingue $r$ de $r-f$; por eso puede usarse la recompensa implícita sin conocer $Z(x)$.
  \item Numéricamente, con $\beta=1$, $\piref=(\tfrac12,\tfrac12)$ y $r'=(6,5)$: $\pi^{*}(0)=e^{6}/(e^{6}+e^{5})=e/(e+1)=0{,}731$.
\end{enumerate}
\resultado \textbf{$\pi^{*}$ y la pérdida de Bradley--Terry son invariantes a sumar $f(x)$; como $\beta\log Z(x)$ es un término de esa forma, desaparece en DPO.}
\comprobacion $0{,}731$ coincide con el valor del \cref{eje:3-beta} sin $f$.
\end{solucion}

% ---------------------------------------------------------------------
\section{Tema 6}

\begin{solucion}{eje:6-cai}
\planteamiento seguir la \cref{fig:cai}: la fase~1 es supervisada; la fase~2, RL con preferencias de IA.
\pasos
\begin{enumerate}[label=\arabic*.]
  \item Fase~1 (SL-CAI): criticar una respuesta según un principio (y revisarla) $\to$ SFT sobre las respuestas revisadas.
  \item Fase~2 (RL-CAI): la IA elige la mejor de dos respuestas $\to$ entrenar el modelo de preferencias $\to$ RL contra él.
\end{enumerate}
\resultado \textbf{criticar (F1) $\to$ SFT sobre revisiones (F1) $\to$ la IA elige (F2) $\to$ modelo de preferencias (F2) $\to$ RL (F2).}
\comprobacion cada paso necesita el resultado del anterior: no hay modelo de preferencias sin pares etiquetados, ni RL sin modelo de preferencias, y los pares de la fase~2 los genera el modelo de la fase~1.
\end{solucion}

\begin{solucion}{eje:6-grpo}
\datos $G=4$; $r=(1,0,0,0)$; desviación poblacional.
\planteamiento \cref{def:grpo-ventaja}.
\pasos
\begin{enumerate}[label=\arabic*.]
  \item Media: $1/4=0{,}25$.
  \item Desviaciones respecto a la media: $0{,}75$ y $-0{,}25$ (tres veces).
  \item Varianza: $(0{,}5625+3\cdot 0{,}0625)/4=0{,}75/4=0{,}1875$; desviación $\sqrt{0{,}1875}=0{,}4330$.
  \item $\widehat A_1=0{,}75/0{,}4330=1{,}732$; $\widehat A_{2,3,4}=-0{,}25/0{,}4330=-0{,}5774$.
  \item Si las cuatro valen~1: media~1, desviación~0; con $\varepsilon_0>0$ todas las ventajas son~0.
\end{enumerate}
\resultado $\mathbf{\widehat A=(1{,}732;\ -0{,}577;\ -0{,}577;\ -0{,}577)}$; con recompensas iguales, $\widehat A=0$ y el grupo no aporta gradiente.
\comprobacion $1{,}732-3\cdot 0{,}577\approx 0$ y $(1{,}732^{2}+3\cdot 0{,}577^{2})/4=(3+1)/4=1$, de acuerdo con la \cref{prop:grpo}.
\end{solucion}

\begin{solucion}{eje:6-drgrpo}
\datos grupos $(1,1,0,0)$ y $(1,0,0,0)$; desviación muestral (entre $G-1=3$).
\planteamiento restar la media del grupo y, en GRPO, dividir entre la desviación.
\pasos
\begin{enumerate}[label=\arabic*.]
  \item Grupo $(1,1,0,0)$: media $0{,}5$; desviaciones $\pm 0{,}5$; varianza muestral $4\cdot 0{,}25/3=1/3$; desviación $0{,}5774$.
  \item GRPO: $\pm 0{,}5/0{,}5774=\pm 0{,}866$. Dr.~GRPO: $\pm 0{,}5$.
  \item Grupo $(1,0,0,0)$: media $0{,}25$; varianza muestral $(0{,}5625+3\cdot 0{,}0625)/3=0{,}25$; desviación $0{,}5$.
  \item GRPO: $0{,}75/0{,}5=1{,}5$ y $-0{,}25/0{,}5=-0{,}5$. Dr.~GRPO: $0{,}75$ y $-0{,}25$.
  \item Factor de amplificación $1/\desv$: $1/0{,}5774=1{,}732$ en el primer grupo y $1/0{,}5=2$ en el segundo.
\end{enumerate}
\resultado \textbf{GRPO: $(0{,}866;\ 0{,}866;\ -0{,}866;\ -0{,}866)$ y $(1{,}5;\ -0{,}5;\ -0{,}5;\ -0{,}5)$; Dr.~GRPO: $(0{,}5;\ 0{,}5;\ -0{,}5;\ -0{,}5)$ y $(0{,}75;\ -0{,}25;\ -0{,}25;\ -0{,}25)$. El grupo $(1,0,0,0)$, de menor desviación, se amplifica más ($\times 2$ frente a $\times 1{,}732$).}
\comprobacion en ambos métodos las ventajas de cada grupo suman~0. Dividir por la desviación da más peso a las instrucciones con poca variación (casi todo aciertos o casi todo fallos), que es la crítica de la \cref{obs:drgrpo}.
\end{solucion}

\begin{solucion}{eje:6-longitud}
\datos longitudes $n_1=2$ y $n_2=4$; longitud máxima~4; $A=3$; $\rho_t=1$; media final sobre dos respuestas.
\planteamiento escribir cada pérdida como suma de $\rho_tA$ dividida por su normalizador, promediar y derivar respecto a cada $\rho_t$.
\pasos
\begin{enumerate}[label=\arabic*.]
  \item Sumas por respuesta: $2\cdot 3=6$ y $4\cdot 3=12$.
  \item \texttt{masked\_\allowbreak{}mean}: $6/2=3$ y $12/4=3$; media $(3+3)/2=3$.
  \item \texttt{masked\_\allowbreak{}normalize}: $6/4=1{,}5$ y $12/4=3$; media $(1{,}5+3)/2=2{,}25$.
  \item Derivada con \texttt{masked\_\allowbreak{}mean}: $A/(2n_i)$; respuesta corta $3/4=0{,}75$, larga $3/8=0{,}375$.
  \item Derivada con \texttt{masked\_\allowbreak{}normalize}: $A/(2\cdot 4)=0{,}375$ en todos los tokens.
\end{enumerate}
\resultado \textbf{\texttt{masked\_\allowbreak{}mean}: valores $(3,3)$ y derivadas $0{,}75$ por token corto y $0{,}375$ por token largo. \texttt{masked\_\allowbreak{}normalize}: valores $(1{,}5;\ 3)$ y derivada $0{,}375$ en todos los tokens.}
\comprobacion con \texttt{masked\_\allowbreak{}mean} cada respuesta pesa lo mismo en total ($2\cdot 0{,}75=4\cdot 0{,}375=1{,}5$), así que un token de una respuesta corta pesa el doble; con ventaja negativa, los tokens de las respuestas largas se penalizan menos, que es el sesgo de longitud de la \cref{obs:drgrpo}.
\end{solucion}

\begin{solucion}{eje:6-cero}
\datos recompensas binarias independientes con probabilidad de acierto $p$; grupos de tamaño $G$.
\planteamiento el grupo no aporta gradiente si todas las recompensas son iguales: todas 1 (probabilidad $p^{G}$) o todas 0 (probabilidad $(1-p)^{G}$).
\pasos
\begin{enumerate}[label=\arabic*.]
  \item $(0{,}5;\ 8)$: $2\cdot 0{,}5^{8}=2/256=0{,}0078$.
  \item $(0{,}9;\ 8)$: $0{,}9^{8}+0{,}1^{8}=0{,}4305+10^{-8}=0{,}4305$.
  \item $(0{,}9;\ 4)$: $0{,}9^{4}+0{,}1^{4}=0{,}6561+0{,}0001=0{,}6562$.
\end{enumerate}
\resultado \textbf{$0{,}78\,\%$, $43{,}0\,\%$ y $65{,}6\,\%$.} Las instrucciones demasiado fáciles o demasiado difíciles desperdician la mayor parte del cómputo; conviene filtrar el conjunto de entrenamiento por dificultad o aumentar $G$.
\comprobacion la expresión es simétrica en $p\leftrightarrow 1-p$ y mínima en $p=0{,}5$; con $p=0{,}1$ y $G=8$ también da $0{,}4305$.
\end{solucion}

\begin{solucion}{eje:6-gao}
\datos $\mathrm{KL}_{\mathrm{bon}}=\log n-(n-1)/n$; $R_{\mathrm{bon}}=d(1-0{,}25d)$; $R_{\mathrm{RL}}=d(1-0{,}5\log d)$; $d=\sqrt{\KL}$.
\planteamiento evaluar la fórmula y aplicar la \cref{prop:gao}.
\pasos
\begin{enumerate}[label=\arabic*.]
  \item $n=4$: $\ln 4-3/4=1{,}386-0{,}75=0{,}636$ nats. $n=16$: $\ln 16-15/16=2{,}773-0{,}9375=1{,}835$ nats.
  \item (b) $d^{*}=\alpha/(2\beta)=1/0{,}5=2$, luego $\mathrm{KL}^{*}=d^{*2}=4$ nats y $R^{*}=\alpha^{2}/(4\beta)=1/1=1$.
  \item Se busca $n$ con $\ln n-(n-1)/n=4$. Como $(n-1)/n\approx 1$ para $n$ grande, $\ln n\approx 5$ y $n\approx e^{5}\approx 148$; la solución numérica es $n\approx 147$.
  \item (c) $d^{*}=e^{\alpha/\beta-1}=e^{2-1}=e=2{,}718$; $R^{*}=\beta d^{*}=0{,}5e=1{,}359$; $\mathrm{KL}^{*}=e^{2}=7{,}39$ nats.
\end{enumerate}
\resultado \textbf{$\mathrm{KL}_{\mathrm{bon}}=0{,}636$ y $1{,}835$ nats; mejor de $n$: $d^{*}=2$, $\mathrm{KL}^{*}=4$ nats, $R^{*}=1$, con $n\approx 147$; RL: $d^{*}=e\approx 2{,}72$ y $R^{*}\approx 1{,}36$.}
\comprobacion $R_{\mathrm{bon}}(2)=2(1-0{,}5)=1$ y $R_{\mathrm{bon}}(4)=4(1-1)=0$: con el doble de la divergencia óptima, la ganancia verdadera desaparece. $R_{\mathrm{RL}}(e)=e(1-0{,}5)=1{,}359$.
\end{solucion}

\begin{solucion}{eje:6-bon}
\datos $U=F(r(y))\sim\mathcal{U}[0,1]$ bajo $\pi_{\mathrm{init}}$; se eligen $n$ muestras independientes y se devuelve la de mayor recompensa.
\planteamiento calcular la probabilidad de que una respuesta concreta sea la elegida y después la esperanza de $\log(\pi_{\mathrm{bon}}/\pi_{\mathrm{init}})$.
\pasos
\begin{enumerate}[label=\arabic*.]
  \item (a) La respuesta elegida está en un entorno de $y$ si una de las $n$ muestras (hay $n$ posibilidades, excluyentes con probabilidad~1) cae allí y las $n-1$ restantes tienen menor recompensa, lo que ocurre con probabilidad $F(r(y))^{n-1}$. Luego $\pi_{\mathrm{bon}}(y)=n\,F(r(y))^{n-1}\pi_{\mathrm{init}}(y)$.
  \item (b) $\log\bigl(\pi_{\mathrm{bon}}/\pi_{\mathrm{init}}\bigr)=\log n+(n-1)\log U$, con $U=F(r(y))$.
  \item Bajo $\pi_{\mathrm{bon}}$, $U$ es el máximo de $n$ uniformes, con densidad $n u^{n-1}$ en $[0,1]$.
  \item $\E[\log U]=\int_0^{1}n u^{n-1}\log u\,du$. Integrando por partes, $\int_0^{1}u^{n-1}\log u\,du=-1/n^{2}$, luego $\E[\log U]=-1/n$.
  \item $\KL=\log n+(n-1)(-1/n)=\log n-(n-1)/n$.
\end{enumerate}
\resultado $\mathbf{\KL(\pi_{\mathrm{bon}}\,\|\,\pi_{\mathrm{init}})=\log n-(n-1)/n}$.
\comprobacion con $n=1$ se obtiene $0$, como debe; con $n=2$, $\ln 2-0{,}5=0{,}193$ nats; la KL crece como $\log n$, de modo que multiplicar $n$ por $e$ añade aproximadamente un nat.
\end{solucion}
